\begin{tikzpicture}[font=\figurefont\small,text=BookInk,
  box/.style={draw=BookBlue,fill=BookBlue!5,line width=.8pt,
    minimum width=16mm,minimum height=9mm,rounded corners=1mm},
  flow/.style={-{Latex[length=2mm]},draw=BookInk,line width=.9pt},
  note/.style={font=\figurefont\footnotesize,align=center}]
  \node[box] (v0) at (0,0) {$v^{(0)}$};
  \node[box] (v1) at (3.5,0) {$v^{(1)}$};
  \node[box] (v2) at (7,0) {$v^{(2)}$};
  \draw[flow] (v0.east) -- node[above,note] {$T^*$}
    node[below,note,text=BookTeal] {旧误差$\times\gamma$} (v1.west);
  \draw[flow] (v1.east) -- node[above,note] {$T^*$}
    node[below,note,text=BookTeal] {旧误差$\times\gamma$} (v2.west);
  \draw[flow] (v2.east) -- (8.75,0) node[right] {$\cdots$};
  \node[note,text=BookAmber] (e0) at (3.5,1.35) {本轮误差\\$e^{(0)}$};
  \node[note,text=BookAmber] (e1) at (7,1.35) {本轮误差\\$e^{(1)}$};
  \draw[flow,BookAmber] (e0.south) -- (v1.north);
  \draw[flow,BookAmber] (e1.south) -- (v2.north);
  \node[note] at (4,-1.3)
    {$E_1\leq\gamma E_0+\delta$\\[2pt]
     $E_2\leq\gamma^2 E_0+\gamma\delta+\delta$};
  \node[note,text=BookMuted] at (4,-2.4)
    {$E_k=\lVert v^{(k)}-V^*\rVert_\infty$，$\lVert e^{(j)}\rVert_\infty\leq\delta$\\
     最终策略提取的贪心误差另计，不重复加入本轮$\delta$};
\end{tikzpicture}
